% Bagian: turing-machines
% Bab: machines-computations
% Subbagian: church-turing-thesis

\documentclass[../../../include/open-logic-section]{subfiles}

\begin{document}

\olfileid[id]{tur}{mac}{ctt}
\olsection{Tesis Church--Turing}

Mesin Turing dimaksudkan sebagai pengganti yang presisi bagi konsep prosedur
efektif. Turing berpendapat bahwa siapa pun yang memahami baik konsep prosedur
efektif maupun konsep mesin Turing akan memiliki intuisi bahwa apa pun yang
dapat dilakukan melalui prosedur efektif juga dapat dilakukan oleh mesin
Turing. Klaim ini didukung oleh fakta bahwa semua pengganti presisi lain yang
diusulkan bagi konsep prosedur efektif ternyata ekuivalen secara ekstensional
dengan konsep mesin Turing---yakni semuanya dapat menghitung tepat himpunan
fungsi yang sama. Klaim ini disebut \emph{tesis Church--Turing}.

\begin{defn}[Tesis Church--Turing]
\emph{Tesis Church--Turing} menyatakan bahwa apa pun yang dapat dihitung
melalui prosedur efektif dapat dihitung oleh mesin Turing.
\end{defn}

Tesis Church--Turing digunakan dalam dua cara. Jenis penggunaan pertama
merupakan alasan untuk bermalas-malasan. Andaikan kita memiliki deskripsi
prosedur efektif untuk menghitung sesuatu, misalnya dalam ``kode semu''.
Kita lalu dapat menggunakan tesis Church--Turing untuk membenarkan klaim bahwa
fungsi yang sama dihitung oleh suatu mesin Turing, sekalipun kita belum
benar-benar membangun mesin tersebut.

Penggunaan lain tesis Church--Turing lebih menarik secara filosofis. Dapat
dibuktikan bahwa ada fungsi-fungsi yang tidak dapat dihitung oleh mesin Turing.
Dari fakta ini dan tesis Church--Turing, dapat disimpulkan bahwa fungsi-fungsi
tersebut tidak dapat dihitung secara efektif dengan prosedur apa pun. Sebab,
jika prosedur semacam itu ada, tesis Church--Turing menyiratkan adanya mesin
Turing untuk prosedur tersebut. Jadi, jika kita dapat membuktikan bahwa tidak
ada mesin Turing yang menghitungnya, tidak mungkin pula ada prosedur efektif.
Secara khusus, tesis Church--Turing digunakan untuk menyatakan bahwa apa yang
disebut masalah penghentian bukan hanya tidak dapat diselesaikan oleh mesin
Turing, melainkan sama sekali tidak dapat diselesaikan secara efektif.

\end{document}
